\documentclass[11pt,compress,t,notes=noshow, xcolor=table]{beamer}

\input{../../style/preamble}
\input{../../latex-math/basic-math}
\input{../../latex-math/basic-ml}
\input{../../latex-math/optim-basics}

\title{Optimization in Machine Learning}

\begin{document}

\titlemeta{
Constrained Optimization
}{
Primal Gradient Descent Methods
}{
figure/primal_gd_projected_step.pdf
}{
\item Why unconstrained GD fails under constraints
\item Projected gradient descent and projected gradients
\item Equality, box, and general convex constraints
\item Convergence facts and ML examples
}

\begin{framev}[fs=small]{Why Constraints Break GD}
\splitVTT[0.46]{
\textbf{Problem}
$$
\min_{\xv \in \S} \fx,
$$
with $\S \subsetneq \R^d$.

\textbf{Vanilla GD step}
$$
\zv^{[t+1]} = \xv^{[t]} - \alpha_t \nabla f(\xv^{[t]}).
$$

\begin{itemize}
\item $\xv^{[t]}$ can be feasible while $\zv^{[t+1]} \notin \S$.
\item Primal methods modify or repair the step so every iterate stays inside $\S$.
\end{itemize}
}{
\imageC[0.95]{figure/primal_gd_projected_step.pdf}
}
\end{framev}


\begin{framev}{Projected Gradient Descent}
\begin{algorithm}[H]
\caption{Projected Gradient Descent}
\begin{algorithmic}[1]
\State Choose a feasible start $\xv^{[0]} \in \S$
\While{stopping criterion not met}
\State $\zv^{[t+1]} \gets \xv^{[t]} - \alpha_t \nabla f(\xv^{[t]})$
\State $\xv^{[t+1]} \gets \Pi_{\S}(\zv^{[t+1]})$
\EndWhile
\end{algorithmic}
\end{algorithm}

\vspace*{-0.2cm}
\begin{itemize}
\item The projection operator is
$$
\Pi_{\S}(\zv) := \argmin_{\xv \in \S} \|\xv - \zv\|_2^2.
$$
\item If the tentative point is already feasible, the projection does nothing.
\item For a closed convex set $\S$, the Euclidean projection is unique.
\end{itemize}
\end{framev}


\begin{framei}[fs=small]{Projected Gradients and Optimality}
\item A constrained optimum $\xvs$ is no longer characterized by $\nabla f(\xvs) = \zero$.
\item Instead, for every $\alpha > 0$,
$$
\xvs = \Pi_{\S}\big(\xvs - \alpha \nabla f(\xvs)\big)
$$
is the fixed-point condition for first-order optimality over a convex set.
\item Define the \textbf{projected gradient mapping}
$$
G_{\alpha}(\xv) := \frac{1}{\alpha}\big(\xv - \Pi_{\S}(\xv - \alpha \nabla f(\xv))\big).
$$
\item $G_{\alpha}(\xv) = \zero$ is the constrained analogue of $\nabla f(\xv) = \zero$.
\item If $\S = \R^d$, then $\Pi_{\S}$ is the identity and $G_{\alpha}(\xv) = \nabla f(\xv)$.
\item Practical stopping rules therefore monitor $\|G_{\alpha}(\xv^{[t]})\|_2$ or the decrease in $f(\xv^{[t]})$.
\end{framei}


\begin{framev}{Equality Constraints: Feasible Directions}
\splitVTT[0.48]{
\textbf{Problem}
$$
\min \fx \qquad \text{s.t. } \qquad \Amat \xv = \bv.
$$

\textbf{Feasible directions}
$$
\Amat \mathbf{d}^{[t]} = \zero.
$$

\begin{itemize}
\item The raw negative gradient usually violates feasibility.
\item Hence we descend only along the tangent / null-space directions.
\end{itemize}
}{
\imageC[0.96]{figure/primal_gd_equality.pdf}
}
\end{framev}


\begin{framei}[fs=small]{Equality Constraints: Algebraic Form}
\item Assume the rows of $\Amat \in \R^{m \times d}$ are linearly independent.
\item The orthogonal projector onto the feasible directions is
$$
\Pi_{\mathrm{null}(\Amat)} = \id - \Amat^\top(\Amat \Amat^\top)^{-1}\Amat.
$$
\item Therefore, the projected-GD update becomes
$$
\xv^{[t+1]} = \xv^{[t]} - \alpha_t \Pi_{\mathrm{null}(\Amat)} \nabla f(\xv^{[t]}).
$$
\item Equivalent viewpoint: parameterize all feasible points as
$$
\xv = \xv_0 + \mathbf{N}\zv, \qquad \Amat \xv_0 = \bv, \qquad \Amat \mathbf{N} = \zero,
$$
and run ordinary GD on the reduced variable $\zv$.
\item Equality constraints are therefore often simpler than general inequalities: feasibility is linear, and the feasible directions have an explicit subspace structure.
\end{framei}


\begin{framev}[fs=small]{Box Constraints: Clipping}
\splitVTT[0.52]{
\textbf{Feasible set}
$$
\S = \left\{\xv \in \R^d \mid
\begin{array}{l}
l_j \le x_j \le u_j \\
j = 1,\dots,d
\end{array}
\right\},
$$

\textbf{Projection}
$$
\left[\Pi_{\S}(\zv)\right]_j
= \min\{u_j, \max\{l_j, z_j\}\}.
$$

\begin{itemize}
\item Each coordinate is clipped independently.
\item No linear system and no QP solve.
\end{itemize}
}{
\imageC[0.95]{figure/primal_gd_box.pdf}
}
\end{framev}


\begin{framei}[fs=small]{General Convex Constraints}
\item For any closed convex feasible set $\S$, projected GD still takes the form
$$
\xv^{[t+1]} = \argmin_{\xv \in \S} \left\| \xv - \big(\xv^{[t]} - \alpha_t \nabla f(\xv^{[t]})\big) \right\|_2^2.
$$
\item The key question is therefore not only how hard it is to differentiate $f$, but also how hard it is to compute $\Pi_{\S}$.
\item Examples with cheap projectors:
\begin{itemize}
\item affine sets and linear equalities
\item box constraints and nonnegativity constraints
\item simplices and some norm balls
\end{itemize}
\item Examples with more expensive projectors:
\begin{itemize}
\item general polyhedra $\{\xv \mid \Amat \xv \le \bv\}$
\item complicated nonlinear convex constraint sets
\end{itemize}
\item This dependence on constraint structure is a recurring theme throughout constrained optimization.
\end{framei}


\begin{framei}[fs=small]{ML Examples for Primal Projected Methods}
\item \textbf{Nonnegative least squares}
$$
\min_{\wv} \|\Xmat \wv - \yv\|_2^2 \qquad \text{s.t. } \qquad \wv \ge \zero
$$
uses the box/nonnegativity projector.
\item \textbf{Probability vectors} in multinomial models satisfy
$$
\thetav \ge \zero, \qquad \one^\top \thetav = 1,
$$
so every step must remain on the simplex.
\item \textbf{Bounded parameter models} and \textbf{budget constraints} naturally yield lower and upper bounds on variables.
\item In all of these cases, the \enquote{primal} viewpoint is intuitive:
\begin{itemize}
\item work directly with the original coefficients
\item keep iterates feasible
\item exploit the geometry of the feasible set
\end{itemize}
\end{framei}


\begin{framei}[fs=small]{Step Sizes and Convergence}
\item Assume $f$ is convex with $L$-Lipschitz continuous gradient and $\S$ is closed and convex.
\item With a constant step size $0 < \alpha_t \le 1 / L$ or with backtracking line search, projected GD keeps all iterates feasible and converges to a minimizer.
\item Typical rates mirror unconstrained GD:
\begin{itemize}
\item convex case: $f(\xv^{[t]}) - f(\xvs) = \order(1/t)$
\item $\mu$-strongly convex case: linear convergence
\end{itemize}
\item The computational cost per iteration is
$$
\text{one gradient evaluation} + \text{one projection}.
$$
\item Hence projected GD is attractive exactly when the projector is simple enough that this extra step is cheap.
\end{framei}


\begin{framev}{Take-Away}
\begin{enumerate}
\item Start from a feasible point.
\item Take a standard gradient step in the original variables.
\item Project back onto the feasible set.
\item Exploit structure aggressively:
\end{enumerate}

\begin{itemize}
\item equalities $\rightarrow$ null-space / tangent-space projection
\item box constraints $\rightarrow$ clipping
\item simple convex sets $\rightarrow$ specialized projectors
\end{itemize}

\vfill
\oneliner{Projected GD is the default primal baseline when projection onto $\S$ is cheap.}

\vfill
Next, we ask what happens when we update only one coordinate at a time.
\end{framev}

\endlecture
\end{document}
